\chapter{Il livello di rete e i router}

\begin{center}
\begin{tikzpicture}[node distance=0pt]
  \node[lay app, minimum width=3.4cm] (a) {Applicazione};
  \node[lay tra, minimum width=3.4cm, below=of a] (t) {Trasporto};
  \node[lay net, minimum width=3.4cm, below=of t, line width=1.5pt, draw=netred] (n) {\bfseries Rete};
  \node[lay link, minimum width=3.4cm, below=of n] (l) {Collegamento};
  \node[lay phy, minimum width=3.4cm, below=of l] (p) {Fisico};
  \draw[-{Stealth}, thick, netred] ([xshift=0.3cm]n.east) -- ++(1,0) node[right, lbl, text=netred] {siamo qui};
\end{tikzpicture}
\end{center}

% ══════════════════════════════════════════════════════════════
\section{Dal trasporto alla rete}
% ══════════════════════════════════════════════════════════════

Il compito del livello di rete è permettere a host remoti di comunicare:

\begin{itemize}
    \item \textbf{sull'host mittente}: prendere i segmenti dal livello di trasporto, incapsulare ciascuno in un
          \textbf{datagramma} e inviare i datagrammi nella rete;
    \item \textbf{sull'host ricevente}: ricevere i datagrammi dalla rete, estrarne i segmenti di trasporto e
          consegnarli al livello di trasporto.
\end{itemize}

A questo livello lavorano alcuni dei protocolli più importanti delle reti (IP, DHCP, NAT, ICMP, i protocolli di
routing) e dispositivi speciali: i \textbf{router}.

\begin{figure}[H]
\centering
\begin{minipage}{0.45\linewidth}\centering
  \includegraphics[width=0.9\linewidth]{cap13/stack_troncato.png}
\end{minipage}\hfill
\begin{minipage}{0.51\linewidth}
All'interno della rete ci sono nodi, i \textbf{router}, che inoltrano i datagrammi ai nodi adiacenti fino all'host
di destinazione.
\begin{itemize}
  \item Nelle figure i router hanno una \textbf{pila troncata} (rete, collegamento, fisico): non eseguono applicazioni
        né protocolli di trasporto.
  \item L'atto di un datagramma che viene reindirizzato da un router si chiama \textbf{hop} (salto).
\end{itemize}
\end{minipage}
\caption{Solo gli host hanno la pila completa; i router si fermano al livello di rete.}
\end{figure}

\begin{figure}[H]
\centering
\begin{tikzpicture}[every node/.style={font=\scriptsize\sffamily}]
  \begin{scope}
    \node[lay app, minimum width=1.8cm] at (0,2.4) {applicazione};
    \node[lay tra, minimum width=1.8cm] at (0,1.8) {trasporto};
    \node[lay net, minimum width=1.8cm] at (0,1.2) {rete};
    \node[lay link, minimum width=1.8cm] at (0,0.6) {collegamento};
    \node[lay phy, minimum width=1.8cm] at (0,0) {fisico};
    \node at (0,-0.6) {\icon[0.7cm]{pc}};
  \end{scope}
  \foreach \x/\n in {3.6/1, 7.2/2} {
    \begin{scope}[xshift=\x cm]
      \node[lay net, minimum width=1.8cm] at (0,1.2) {rete};
      \node[lay link, minimum width=1.8cm] at (0,0.6) {collegamento};
      \node[lay phy, minimum width=1.8cm] at (0,0) {fisico};
      \node at (0,-0.6) {\icon[0.9cm]{router}};
    \end{scope}
  }
  \begin{scope}[xshift=10.8cm]
    \node[lay app, minimum width=1.8cm] at (0,2.4) {applicazione};
    \node[lay tra, minimum width=1.8cm] at (0,1.8) {trasporto};
    \node[lay net, minimum width=1.8cm] at (0,1.2) {rete};
    \node[lay link, minimum width=1.8cm] at (0,0.6) {collegamento};
    \node[lay phy, minimum width=1.8cm] at (0,0) {fisico};
    \node at (0,-0.6) {\icon[0.4cm]{server}};
  \end{scope}
  \draw[-{Stealth}, very thick, netred] (0,2.1) -- (0,-0.25) -- (2.6,-0.25) -- (2.6,1.2) -- (4.6,1.2) -- (4.6,-0.25) -- (6.2,-0.25) -- (6.2,1.2) -- (8.2,1.2) -- (8.2,-0.25) -- (9.8,-0.25) -- (9.8,2.1);
  \node[text=netred] at (1.8,-1.1) {hop 1}; \node[text=netred] at (5.4,-1.1) {hop 2}; \node[text=netred] at (9,-1.1) {hop 3};
\end{tikzpicture}
\caption{Il datagramma sale fino al livello di rete in ogni router, che decide dove inoltrarlo, e riscende.}
\end{figure}

% ══════════════════════════════════════════════════════════════
\section{Che servizi offre il livello di rete}
% ══════════════════════════════════════════════════════════════

Il livello di rete è responsabile della \textbf{consegna host-a-host}. In teoria, insieme alla consegna potrebbe
offrire servizi come:

\begin{itemize}
    \item \textbf{consegna garantita}: un pacchetto inviato arriverà prima o poi a destinazione;
    \item \textbf{consegna garantita con ritardo limitato}: arriverà entro un certo tempo (per esempio 100 ms);
    \item \textbf{consegna in ordine}: i pacchetti arrivano nell'ordine di invio;
    \item \textbf{banda minima garantita}: se il mittente non supera una certa velocità (per esempio 1 Mb/s), tutti i
          pacchetti arrivano;
    \item \textbf{sicurezza}: cifratura e decifratura di tutti i datagrammi.
\end{itemize}

In realtà, \textbf{nessuno} di questi servizi è offerto in generale dalle reti.

\dfn{Servizio best effort}{
    Il livello di rete di Internet offre un solo servizio, detto \textbf{best effort} (``del proprio meglio''): la
    rete \textbf{fa del suo meglio} per consegnare un pacchetto dalla sorgente alla destinazione, ma non garantisce
    né che arrivi, né che arrivi in ordine, né ritardi massimi, né banda minima.
}

Nonostante la sua semplicità, il modello best effort, insieme a una buona disponibilità di banda, si è
dimostrato \textbf{adeguato}: Netflix, le videochiamate, le conferenze in tempo reale funzionano tutte così.

% ══════════════════════════════════════════════════════════════
\section{I router nella rete}
% ══════════════════════════════════════════════════════════════

I router sono nodi speciali con \textbf{più collegamenti in ingresso e in uscita}. Offrono due funzioni di rete
fondamentali.

\dfn{Forwarding e routing}{
    \begin{itemize}
      \item \textbf{Forwarding} (inoltro): quando un pacchetto arriva su un collegamento di ingresso, il router lo
            sposta sul \textbf{collegamento di uscita} appropriato. È un'azione \textbf{locale} al router.
      \item \textbf{Routing} (instradamento): decidere il \textbf{percorso} che i pacchetti devono seguire dalla
            sorgente alla destinazione. È un processo che riguarda \textbf{tutta la rete}; gli algoritmi che
            calcolano i percorsi sono gli \textbf{algoritmi di routing}.
    \end{itemize}
}

Durante il forwarding un router può anche:

\begin{enumerate}
    \item \textbf{bloccare} un pacchetto (se proviene da un host noto come malevolo o è diretto a un host vietato);
    \item \textbf{scartare} un pacchetto scaduto;
    \item \textbf{duplicare} un pacchetto e inviarlo su più collegamenti;
    \item \textbf{modificare} un pacchetto (per esempio per segnalare congestione).
\end{enumerate}

\info{L'analogia del viaggio in auto}{
    Il \textbf{routing} è pianificare il viaggio da Napoli a Milano, scegliendo le autostrade da percorrere. Il
    \textbf{forwarding} è quello che si fa a ogni svincolo: guardare il cartello e imboccare l'uscita giusta. Il primo
    richiede una visione d'insieme e si fa una volta; il secondo è rapido, locale e si ripete a ogni incrocio.
}

\begin{figure}[H]
\centering
\begin{tikzpicture}[every node/.style={font=\footnotesize\sffamily}]
  \begin{scope}
    \node[font=\small\bfseries\sffamily] at (1.5,1.6) {Forwarding};
    \node (r) at (1.5,0) {\icon[1.3cm]{router}};
    \draw[msg, very thick, netblue] (-0.8,0.5) -- node[above] {P} (0.8,0.1);
    \draw[msg, very thick, netgreen] (2.2,0.1) -- (3.8,0.9); \draw[thick, black!30] (2.2,-0.1) -- (3.8,-0.8); \draw[thick, black!30] (0.8,-0.2) -- (-0.8,-0.8);
    \node[below=4pt of r, text width=3cm, align=center] {sposta il pacchetto dall'ingresso all'uscita giusta};
  \end{scope}
  \begin{scope}[xshift=7cm]
    \node[font=\small\bfseries\sffamily] at (2.5,1.6) {Routing};
    \foreach \x/\y [count=\i] in {0/0, 1.5/0.9, 1.6/-0.9, 3.2/0.6, 3.3/-0.7, 4.8/0} \node (n\i) at (\x,\y) {\icon[0.7cm]{router}};
    \foreach \a/\b in {1/2, 1/3, 2/4, 3/5, 2/5, 4/6, 5/6, 3/4} \draw[thick, black!30] (n\a) -- (n\b);
    \draw[pathhl, -{Stealth}] (n1) -- (n2) -- (n4) -- (n6);
    \node[text width=4.6cm, align=center] at (2.4,-1.8) {calcola il percorso migliore attraverso tutta la rete};
  \end{scope}
\end{tikzpicture}
\caption{Forwarding (locale) e routing (globale).}
\end{figure}

% ══════════════════════════════════════════════════════════════
\section{Tipi di router}
% ══════════════════════════════════════════════════════════════

\begin{figure}[H]
\centering
\begin{minipage}{0.4\linewidth}\centering
  \includegraphics[width=0.95\linewidth]{cap13/router_foto.png}
\end{minipage}\hfill
\begin{minipage}{0.56\linewidth}
I router possono essere molto diversi a seconda di uso e tecnologia: per casa o per ufficio, con o senza fili.
Una distinzione importante è quella tra:
\begin{itemize}
  \item \textbf{edge router} (router di bordo): distribuiscono pacchetti \textbf{tra reti diverse}, per esempio
        collegando una rete con l'ISP. È il tipo più comune;
  \item \textbf{core router} (router di nucleo): distribuiscono pacchetti \textbf{all'interno di una stessa rete}.
        Sono usati anche sulla dorsale (\emph{backbone}) di Internet, dove devono smaltire trasferimenti di dati enormi.
\end{itemize}
\end{minipage}
\caption{Router per uso domestico e aziendale (in alto), un edge router e un core router da dorsale (in basso).}
\end{figure}

% ══════════════════════════════════════════════════════════════
\section{Data plane e control plane}
% ══════════════════════════════════════════════════════════════

\dfn{I due piani di un router}{
    \begin{itemize}
      \item Il \textbf{data plane} (piano dei dati, o di forwarding) comprende l'azione locale di trasferire un
            pacchetto da un'interfaccia di ingresso a quella di uscita appropriata. È un'operazione \textbf{veloce}
            (pochi nanosecondi), di solito realizzata in \textbf{hardware}.
      \item Il \textbf{control plane} (piano di controllo) comprende il processo, esteso a tutta la rete, che
            determina i percorsi da un estremo all'altro. È più \textbf{lento} (secondi), e spesso realizzato in
            \textbf{software}.
    \end{itemize}
}

Nel data plane c'è la \textbf{forwarding table} (tabella di inoltro): una tabella che associa gli indirizzi di
destinazione ai collegamenti di uscita. Il router inoltra un pacchetto esaminando uno o più campi del suo
header; il valore trovato nella tabella indica l'interfaccia su cui inoltrarlo.

\warning{La tabella non contiene i singoli indirizzi}{
    Una forwarding table \textbf{non} elenca gli indirizzi dei singoli host (sarebbero miliardi): associa alle porte
    di uscita \textbf{gruppi di indirizzi}, cioè intervalli individuati da un prefisso (le \emph{subnet}).
}

\subsection{Centralizzato o decentralizzato}

Poiché prima di arrivare a destinazione si attraversano molti router, il contenuto di una forwarding table va
determinato raccogliendo informazioni da router diversi. Lo si fa con una o più \textbf{tabelle di routing}
(control plane), da cui si ricava la forwarding table (data plane). Ci sono due modi:

\begin{figure}[H]
\centering
\begin{tikzpicture}[every node/.style={font=\scriptsize\sffamily}]
  \begin{scope}
    \node[font=\footnotesize\bfseries\sffamily] at (2.3,2.7) {Decentralizzato};
    \foreach \x [count=\i] in {0,1.5,3,4.5} {
      \node[draw, fill=lnet!40, minimum width=1.1cm, minimum height=0.5cm] (cp\i) at (\x,1.7) {routing};
      \node[draw, fill=llink!60, minimum width=1.1cm, minimum height=0.5cm] (dp\i) at (\x,0.9) {tabella};
      \draw[-{Stealth}] (cp\i) -- (dp\i);
      \node at (\x,0.1) {\icon[0.6cm]{router}};
    }
    \foreach \i/\j in {1/2, 2/3, 3/4} \draw[{Stealth}-{Stealth}, netred, thick] (cp\i) -- (cp\j);
    \node[text width=4.8cm, align=center] at (2.25,-0.7) {ogni router ha un componente di routing che dialoga con quelli degli altri};
  \end{scope}
  \begin{scope}[xshift=8cm]
    \node[font=\footnotesize\bfseries\sffamily] at (2.3,2.7) {Centralizzato (SDN)};
    \node[draw, fill=lnet!60, rounded corners, minimum width=4.8cm, minimum height=0.6cm] (ctrl) at (2.25,1.9) {controller remoto (software)};
    \foreach \x [count=\i] in {0,1.5,3,4.5} {
      \node[draw, fill=llink!60, minimum width=1.1cm, minimum height=0.5cm] (dp\i) at (\x,0.9) {tabella};
      \draw[-{Stealth}, netred, thick] (ctrl.south -| dp\i) -- (dp\i);
      \node at (\x,0.1) {\icon[0.6cm]{router}};
    }
    \node[text width=4.8cm, align=center] at (2.25,-0.7) {un controller separato calcola e distribuisce le tabelle di inoltro};
  \end{scope}
\end{tikzpicture}
\caption{Control plane decentralizzato e centralizzato.}
\end{figure}

\begin{itemize}
    \item \textbf{Decentralizzato} (distribuito): ogni router ha un componente di routing che comunica con i componenti
          di routing degli altri router. È stato a lungo l'approccio dominante.
    \item \textbf{Centralizzato}: un \textbf{controller remoto}, fisicamente separato dai router, calcola le tabelle e
          le distribuisce. Può trovarsi in un data center remoto, affidabile e ridondato, gestito dall'ISP o da terzi.
\end{itemize}

\info{Software-Defined Networking}{
    L'approccio centralizzato è alla base del \textbf{Software-Defined Networking} (SDN): la rete è ``definita dal
    software'' perché il controller che calcola le tabelle e dialoga con i router è un software centralizzato,
    talvolta anche open source.
}

% ══════════════════════════════════════════════════════════════
\section{Le componenti di un router}
% ══════════════════════════════════════════════════════════════

\begin{figure}[H]
\centering
\begin{tikzpicture}[every node/.style={font=\footnotesize\sffamily}]
  \draw[thick, rounded corners, fill=black!3] (-0.6,-3.3) rectangle (9.6,1.5);
  \node[anchor=north west, font=\small\bfseries\sffamily] at (-0.5,1.45) {Router};
  \foreach \y/\i in {0/1, -1.2/2, -2.6/N} {
    \node[draw, fill=lnet!40, minimum width=2cm, minimum height=0.8cm, align=center] (in\i) at (1,\y) {porta di\\ingresso \i};
    \node[draw, fill=llink!60, minimum width=2cm, minimum height=0.8cm, align=center] (out\i) at (8,\y) {porta di\\uscita \i};
    \draw[-{Stealth}, very thick] (-1.6,\y) -- (in\i.west);
    \draw[-{Stealth}, very thick] (out\i.east) -- (10.6,\y);
  }
  \node at (1,-1.9) {$\vdots$}; \node at (8,-1.9) {$\vdots$};
  \node[draw, fill=ltra!50, minimum width=2.4cm, minimum height=3.4cm, align=center] (sf) at (4.5,-1.3) {switching\\fabric};
  \foreach \i in {1,2,N} { \draw[-{Stealth}, thick] (in\i.east) -- (in\i.east -| sf.west); \draw[-{Stealth}, thick] (out\i.west -| sf.east) -- (out\i.west); }
  \node[draw, fill=lapp!50, rounded corners, minimum width=3cm, align=center] (rp) at (4.5,2.4) {processore di routing\\(control plane)};
  \draw[{Stealth}-{Stealth}, dashed, thick, netred] (rp.south) -- node[right, text=netred] {tabelle} (sf.north);
  \draw[-{Stealth}, dashed, netred] (rp.west) -| (in1.north);
  \node[anchor=east] at (-1.7,-1.3) {collegamenti di ingresso};
  \node[anchor=west, text width=2cm] at (10.7,-1.3) {collegamenti di uscita};
  \node[text=black!60] at (4.5,-3.0) {data plane (hardware)};
\end{tikzpicture}
\caption{Architettura di un router.}
\end{figure}

\begin{itemize}
    \item \textbf{Porte di ingresso} (da non confondere con le porte del trasporto): sono le interfacce fisiche
          collegate ai collegamenti in ingresso. Il loro compito è:
    \begin{itemize}
      \item consultare la forwarding table (\emph{lookup}) e preparare la switching fabric verso la porta di uscita scelta;
      \item inoltrare al processore di routing i pacchetti di controllo (per esempio quelli con le informazioni di routing).
    \end{itemize}
    Il numero di porte va da decine a centinaia (l'edge router Juniper MX2020 supporta fino a 960 porte Ethernet da 10 Gb/s).
    \item \textbf{Switching fabric} (matrice di commutazione): collega le porte di ingresso a quelle di uscita.
    \item \textbf{Porte di uscita}: conservano i pacchetti ricevuti dalla switching fabric e li trasmettono sul
          collegamento in uscita. Le porte possono essere bidirezionali (ingresso e uscita accoppiati).
    \item \textbf{Processore di routing}: esegue i protocolli di routing, mantiene informazioni sullo stato dei
          collegamenti e calcola e aggiorna la forwarding table.
\end{itemize}

% ══════════════════════════════════════════════════════════════
\section{Il forwarding e la longest prefix matching}
% ══════════════════════════════════════════════════════════════

La forwarding table associa \textbf{indirizzi IP} a \textbf{porte di uscita}, così che i pacchetti possano essere
inoltrati sul collegamento giusto verso il nodo successivo (dove, eventualmente, verranno inoltrati di nuovo).

Un indirizzo IP è un numero di \textbf{4 byte} (32 bit), che di solito si scrive in notazione decimale ma che si può
vedere in binario:

\begin{center}
\begin{tabular}{ll}
decimale puntata & \texttt{192.168.1.1} \\
binario & \texttt{11000000 10101000 00000001 00000001} \\
\end{tabular}
\end{center}

Quando un pacchetto arriva su una porta di ingresso si esegue un \textbf{lookup}. Ogni porta ha una
\textbf{copia} della forwarding table, ricevuta dal processore di routing su un bus dedicato (per esempio PCI), per
evitare il collo di bottiglia di interrogare il processore per ogni pacchetto. Più indirizzi possono essere
associati alla stessa porta.

\begin{table}[H]
\centering
\begin{tabular}{llc}
\toprule
\textbf{Prefisso degli indirizzi} & \textbf{Intervallo di indirizzi corrispondente} & \textbf{Interfaccia} \\
\midrule
\texttt{11001000 00010111 00010*** ********} & da \texttt{...00010000 00000000} a \texttt{...00010111 11111111} & 0 \\
\texttt{11001000 00010111 00011000 ********} & da \texttt{...00011000 00000000} a \texttt{...00011000 11111111} & 1 \\
\texttt{11001000 00010111 00011*** ********} & da \texttt{...00011000 00000000} a \texttt{...00011111 11111111} & 2 \\
altrimenti & & 3 \\
\bottomrule
\end{tabular}
\caption{Una forwarding table a 4 porte: ogni riga è un prefisso, cioè un intervallo di indirizzi (i primi 16 bit sono sempre \texttt{11001000 00010111}).}
\end{table}

\warning{Righe che si sovrappongono}{
    Le righe non sono necessariamente mutuamente esclusive. Per esempio l'indirizzo
    \texttt{11001000 00010111 00011000 10101010} corrisponde sia alla riga dell'interfaccia 1 sia a quella
    dell'interfaccia 2. A quale porta va?
}

\dfn{Longest prefix matching}{
    Se un indirizzo corrisponde a più righe della forwarding table, il router lo inoltra verso la porta associata alla
    riga con il \textbf{prefisso più lungo} (\emph{longest prefix matching rule}).
}

\begin{figure}[H]
\centering
\begin{tikzpicture}[every node/.style={font=\small\ttfamily}]
  \node[anchor=west] (addr) at (0,0) {\textcolor{netred}{11001000 00010111 00011000} 10101010};
  \node[anchor=west, font=\small\sffamily] at (8,0) {indirizzo di destinazione};
  \draw[decorate, decoration={brace, amplitude=5pt}, netgreen!60!black, thick] (0.1,0.35) -- (6.05,0.35) node[midway, above=6pt, font=\footnotesize\sffamily, text=netgreen!50!black] {match con la riga dell'interfaccia 1: 24 bit};
  \draw[decorate, decoration={brace, amplitude=5pt, mirror}, netblue, thick] (0.1,-0.35) -- (5.35,-0.35) node[midway, below=6pt, font=\footnotesize\sffamily, text=netblue] {match con la riga dell'interfaccia 2: 21 bit};
  \node[font=\small\bfseries\sffamily, text=netgreen!50!black] at (10,-1.1) {vince l'interfaccia 1!};
\end{tikzpicture}
\caption{Longest prefix matching: tra le righe che corrispondono si sceglie quella con il prefisso più lungo.}
\end{figure}

\begin{esempio}{Dove va il pacchetto?}
Con la tabella precedente:
\begin{itemize}
    \item \texttt{11001000 00010111 00010110 10100001}: corrisponde solo al prefisso \texttt{...00010} $\Rightarrow$ interfaccia \textbf{0};
    \item \texttt{11001000 00010111 00011000 10101010}: corrisponde a \texttt{...00011000} (24 bit) e a \texttt{...00011} (21 bit)
          $\Rightarrow$ interfaccia \textbf{1};
    \item \texttt{11001000 00010111 00011001 00000001}: corrisponde solo a \texttt{...00011} $\Rightarrow$ interfaccia \textbf{2};
    \item \texttt{11001000 00010111 00100000 00000001}: nessun prefisso $\Rightarrow$ interfaccia \textbf{3}.
\end{itemize}
\end{esempio}

\info{Perché è utile}{
    La regola del prefisso più lungo permette di assegnare un \textbf{sotto-spazio} di un blocco di indirizzi a una porta
    di uscita diversa, \textbf{senza ambiguità}: la riga più generica copre il blocco, quella più specifica ``ritaglia''
    un'eccezione. Lo useremo con l'aggregazione degli indirizzi.
}

Il lookup è di solito eseguito \textbf{in hardware}, per essere il più veloce possibile (nanosecondi per le
trasmissioni a gigabit), con memorie dedicate e algoritmi di ricerca avanzati. Determinata la porta di uscita, il
pacchetto entra nella switching fabric; se altre porte stanno usando la fabric, può dover attendere in coda.

\subsection{Match-plus-action}

L'operazione in due passi ``cerca l'indirizzo di destinazione'' (\emph{match}) e ``inoltra'' (\emph{action}) si
chiama \textbf{match-plus-action} ed è eseguita in molti dispositivi di rete:

\begin{itemize}
    \item \textbf{switch}: azione simile a quella dei router (ma su indirizzi di livello 2);
    \item \textbf{firewall}: l'azione è filtrare, cioè bloccare, certi pacchetti in arrivo;
    \item \textbf{NAT} (\emph{Network Address Translator}): l'azione è riscrivere indirizzo e porta di certi pacchetti prima di inoltrarli.
\end{itemize}

% ══════════════════════════════════════════════════════════════
\section{La switching fabric}
% ══════════════════════════════════════════════════════════════

\begin{figure}[H]
\centering
\begin{tikzpicture}[every node/.style={font=\scriptsize\sffamily}]
  \begin{scope}
    \node[font=\footnotesize\bfseries\sffamily] at (1.6,2.4) {1. Memoria};
    \foreach \y [count=\k] in {1.5,0.5,-0.5} { \node[draw, fill=lnet!40, minimum width=0.8cm] (i\k) at (0,\y) {in}; \node[draw, fill=llink!60, minimum width=0.8cm] (o\k) at (3.2,\y) {out}; }
    \node[draw, fill=ltra!50, minimum width=1.2cm, minimum height=2.6cm, align=center] (mem) at (1.6,0.5) {CPU\\+\\memoria};
    \draw[-{Stealth}, netred, thick] (i1) -- (mem); \draw[-{Stealth}, netred, thick] (mem) -- (o3);
    \node[text width=3.6cm, align=center] at (1.6,-1.5) {il processore copia il pacchetto: lento, tipico dei primi router};
  \end{scope}
  \begin{scope}[xshift=5.6cm]
    \node[font=\footnotesize\bfseries\sffamily] at (1.6,2.4) {2. Bus};
    \foreach \y in {1.5,0.5,-0.5} { \node[draw, fill=lnet!40, minimum width=0.8cm] at (0,\y) {in}; \node[draw, fill=llink!60, minimum width=0.8cm] at (3.2,\y) {out}; \draw (0.4,\y) -- (2.8,\y); }
    \draw[line width=4pt, netblue] (1.6,1.9) -- (1.6,-0.9);
    \draw[-{Stealth}, netred, very thick] (0.4,1.5) -- (1.6,1.5) -- (1.6,-0.5) -- (2.8,-0.5);
    \node[text width=3.6cm, align=center] at (1.6,-1.5) {bus condiviso: una porta alla volta, basta per reti piccole};
  \end{scope}
  \begin{scope}[xshift=11.2cm]
    \node[font=\footnotesize\bfseries\sffamily] at (1.6,2.4) {3. Rete di interconnessione};
    \foreach \y in {1.5,0.5,-0.5} { \node[draw, fill=lnet!40, minimum width=0.8cm] at (0,\y) {in}; \draw (0.4,\y) -- (2.6,\y); }
    \foreach \x in {1.0,1.6,2.2} { \draw (\x,1.9) -- (\x,-0.9); \node[draw, fill=llink!60, minimum width=0.5cm, rotate=90] at (\x,-1.1) {\tiny out}; }
    \foreach \x in {1.0,1.6,2.2} \foreach \y in {1.5,0.5,-0.5} \fill[black!40] (\x,\y) circle (1.5pt);
    \draw[netred, very thick, -{Stealth}] (0.4,1.5) -- (2.2,1.5) -- (2.2,-0.8);
    \draw[netgreen!60!black, very thick, -{Stealth}] (0.4,-0.5) -- (1.0,-0.5) -- (1.0,-0.8);
    \node[text width=3.6cm, align=center] at (1.6,-1.9) {punti di incrocio che si aprono e chiudono: più pacchetti in parallelo};
  \end{scope}
\end{tikzpicture}
\caption{I tre modi di realizzare la commutazione.}
\end{figure}

\begin{enumerate}
    \item \textbf{Commutazione tramite memoria}: le porte sono trattate come dispositivi di I/O che scrivono i pacchetti in
          memoria; il processore di routing li copia sulla porta di uscita indicata dalla tabella. È un po' lenta (serve
          accedere alla memoria) ed era comune nei primi router, che erano normali computer.
    \item \textbf{Commutazione tramite bus}: la porta di ingresso trasferisce il pacchetto direttamente alla porta di uscita
          su un bus condiviso, senza intervento del processore. Si serve una porta alla volta, ma spesso basta per router di
          reti locali piccole.
    \item \textbf{Commutazione tramite rete di interconnessione} (\emph{crossbar}): le porte sono collegate da una rete con
          punti di incrocio che possono essere aperti o chiusi per dirigere i pacchetti. Più pacchetti possono essere
          inoltrati \textbf{in parallelo}: è il metodo di molti router moderni.
\end{enumerate}

% ══════════════════════════════════════════════════════════════
\section{Le code sulle porte}
% ══════════════════════════════════════════════════════════════

Poiché la commutazione richiede tempo, porte di ingresso e di uscita hanno \textbf{code} (buffer) in cui i pacchetti
aspettano, come le auto a un semaforo. L'entità delle code non è fissa: dipende dal carico di traffico, dalla velocità
della switching fabric e da quella delle linee. In generale i pacchetti possono arrivare o partire più velocemente o
più lentamente di quanto avvenga la commutazione, e accumularsi in un buffer di ingresso o di uscita, che può
\textbf{traboccare}.

\begin{esempio}{Dove si formano le code}
Un router ha $N$ porte di ingresso e $N$ di uscita; tutte le linee hanno velocità $R_{line}$ pacchetti al secondo e tutte
le porte ricevono pacchetti contemporaneamente. Nel caso peggiore \textbf{tutti} i pacchetti vanno alla \textbf{stessa}
porta di uscita. Detta $R_{switch}$ la velocità della switching fabric:
\begin{itemize}
    \item se $R_{switch} \approx N \cdot R_{line}$: la fabric smaltisce tutto in tempo, le code in \textbf{ingresso} sono
          trascurabili, ma quella in \textbf{uscita} cresce, perché arrivano $N$ volte più pacchetti di quanti il collegamento ne trasmetta;
    \item se $R_{line} < R_{switch} < N \cdot R_{line}$: si formano code \textbf{su entrambe} le porte (l'ingresso aspetta la
          fabric, l'uscita aspetta il collegamento);
    \item se $R_{switch} \approx R_{line}$: la coda in uscita è trascurabile, ma quelle in \textbf{ingresso} sono consistenti.
\end{itemize}
\end{esempio}

\warning{Head-of-the-line blocking}{
    Quando si accodano pacchetti in ingresso può capitare che un pacchetto resti bloccato perché quello \textbf{davanti
    a lui} nella stessa coda aspetta una porta di uscita occupata, anche se la sua porta di uscita sarebbe libera. È il
    \emph{blocco in testa alla coda} (HOL blocking), uno dei motivi per cui le code in ingresso vanno evitate.
}

% ══════════════════════════════════════════════════════════════
\section{La schedulazione dei pacchetti}
% ══════════════════════════════════════════════════════════════

È normale che più pacchetti, anche da porte di ingresso diverse, debbano uscire dalla stessa porta. L'accesso dei
pacchetti in coda al collegamento di uscita va quindi \textbf{schedulato}. I tre approcci più noti sono:

\begin{itemize}
    \item \textbf{FIFO} (\emph{First-In-First-Out}, o FCFS, \emph{First-Come-First-Served}): basato sull'ordine di arrivo;
    \item \textbf{a priorità} (\emph{priority queuing}): basato sull'importanza dei pacchetti;
    \item \textbf{round-robin}: i pacchetti sono divisi in classi, servite a turno.
\end{itemize}

\begin{figure}[H]
\centering
\begin{tikzpicture}[every node/.style={font=\scriptsize\sffamily}]
  \begin{scope}
    \node[font=\footnotesize\bfseries\sffamily] at (1.3,2.2) {FIFO};
    \draw[thick] (0,0) -- (0,1.5) -- (2.6,1.5); \draw[thick] (0,0) -- (2.6,0); \draw[thick] (2.6,0) -- (2.6,1.5);
    \foreach \x/\n in {0.3/3, 1.1/2, 1.9/1} { \fill[netblue!50] (\x,0.4) rectangle ++(0.5,0.7); \node at (\x+0.25,0.75) {\n}; }
    \draw[-{Stealth}, very thick] (-1,0.75) -- (0,0.75); \draw[-{Stealth}, very thick] (2.6,0.75) -- (3.6,0.75) node[right] {collegamento};
    \node[text width=3.4cm, align=center] at (1.3,-0.7) {esce per primo chi è arrivato per primo; a coda piena si scarta l'ultimo arrivato (\emph{drop-tail})};
  \end{scope}
  \begin{scope}[xshift=5.8cm]
    \node[font=\footnotesize\bfseries\sffamily] at (1.3,2.2) {Priorità};
    \foreach \y/\c/\t in {1.0/netred/alta, 0/netblue/bassa} {
      \draw[thick] (0,\y) rectangle (2.2,\y+0.7);
      \foreach \x in {1.0,1.6} \fill[\c!50] (\x,\y+0.1) rectangle ++(0.45,0.5);
      \node[anchor=east] at (-0.1,\y+0.35) {\t};
      \draw[-{Stealth}, thick] (2.2,\y+0.35) -- (2.9,0.85);
    }
    \node[circle, draw, fill=white] (sw) at (3.1,0.85) {};
    \draw[-{Stealth}, very thick] (sw) -- (4,0.85);
    \node[text width=3.6cm, align=center] at (1.6,-0.7) {si serve sempre la classe più importante non vuota; dentro la classe si va in FIFO};
  \end{scope}
  \begin{scope}[xshift=11.4cm]
    \node[font=\footnotesize\bfseries\sffamily] at (1.3,2.2) {Round-robin / WFQ};
    \foreach \y/\c/\t in {1.0/netgreen/{classe 1 (peso 2)}, 0/netorange/{classe 2 (peso 1)}} {
      \draw[thick] (0,\y) rectangle (2.2,\y+0.7);
      \foreach \x in {1.0,1.6} \fill[\c!50] (\x,\y+0.1) rectangle ++(0.45,0.5);
      \node[anchor=south] at (1.1,\y+0.7) {\t};
      \draw[-{Stealth}, thick] (2.2,\y+0.35) -- (2.9,0.85);
    }
    \node[circle, draw, fill=white] (sw) at (3.1,0.85) {};
    \draw[-{Stealth}, very thick] (sw) -- (4,0.85);
    \draw[-{Stealth}, netred] (3.1,1.35) arc (90:-200:0.5);
    \node[text width=3.6cm, align=center] at (1.6,-0.7) {le classi si alternano; con WFQ ogni classe ha un peso che ne fissa la quota};
  \end{scope}
\end{tikzpicture}
\caption{Tre politiche di schedulazione della coda di uscita.}
\end{figure}

\subsection{FIFO (First-Come-First-Served)}

\begin{itemize}
    \item Se il collegamento di uscita è occupato, i pacchetti in arrivo alla porta di uscita vengono \textbf{bufferizzati}.
    \item Se non c'è spazio sufficiente, serve una politica di scarto: la più comune è scartare il pacchetto appena arrivato
          (\textbf{drop-tail}); in approcci più sofisticati si possono rimuovere anche pacchetti già in coda per fare posto ai nuovi.
    \item Un pacchetto viene tolto dalla coda solo quando è stato \textbf{completamente trasmesso}.
    \item I pacchetti vengono trasmessi nello stesso ordine in cui sono arrivati.
\end{itemize}

\subsection{Priority queuing}

\begin{itemize}
    \item All'arrivo in coda i pacchetti vengono \textbf{classificati} in classi di priorità (per esempio in base alle porte TCP/UDP).
    \item Un operatore può configurare la coda perché alcuni pacchetti (informazioni di gestione della rete, voce su IP in
          tempo reale, \dots) abbiano priorità sul traffico degli utenti o non in tempo reale.
    \item Ogni classe ha la propria coda: si trasmettono per primi i pacchetti della classe \textbf{più prioritaria non vuota};
          dentro una stessa classe si procede in FIFO.
\end{itemize}

\subsection{Round-robin queuing}

\begin{itemize}
    \item I pacchetti sono ancora divisi in classi, ma le classi vengono \textbf{alternate} invece che scelte per priorità.
    \item Un'implementazione comune è il \textbf{Weighted Fair Queuing} (WFQ): i pacchetti vengono classificati e accodati
          nella propria area di attesa, e a ogni classe è associato un \textbf{peso} che determina con quale frequenza viene
          servita rispetto alle altre.
\end{itemize}

% ══════════════════════════════════════════════════════════════
\section{Riepilogo}
% ══════════════════════════════════════════════════════════════

\riepilogo{
\begin{itemize}
    \item Il livello di rete consegna datagrammi \textbf{da host a host}, con servizio \textbf{best effort}.
    \item I router fanno \textbf{forwarding} (locale, data plane, hardware, nanosecondi) e \textbf{routing} (globale, control
          plane, software, secondi); il control plane può essere distribuito o centralizzato (SDN).
    \item Un router ha porte di ingresso, switching fabric (memoria, bus, crossbar), porte di uscita e processore di routing.
    \item La forwarding table associa \textbf{prefissi} a interfacce; vince il \textbf{prefisso più lungo}.
    \item Le code possono formarsi in ingresso e in uscita; la schedulazione può essere FIFO, a priorità o round-robin (WFQ).
\end{itemize}
}
